\section{Background}


\subsection{Hanzo Base Architecture}

Hanzo Base runs one SQLite database per organization. Each database is
encrypted at rest with sqlcipher (AES-256-CBC, PBKDF2-derived key). The
encryption key is stored in Hanzo KMS (HIP-0027) and fetched at startup.
Base exposes a REST API, a realtime SSE subscription layer, and a JavaScript
plugin runtime (Goja VM). All data operations go through SQLite's WAL
(Write-Ahead Log) mode.

WAL mode is critical to this proposal. In WAL mode, SQLite writes all
changes to a separate WAL file before checkpointing them into the main
database. This creates a sequential, append-only stream of changes---exactly
the primitive that continuous replication requires.

\subsection{Litestream}

Litestream~\cite{litestream} is an open-source tool by Ben Johnson that
continuously replicates SQLite databases to S3-compatible storage. It works
by monitoring the SQLite WAL file and streaming new frames to a replica as
they are written. On restore, it downloads the most recent snapshot and
replays WAL frames to reconstruct the database to its last replicated state.

Litestream operates as a separate process alongside the application. It
requires no application code changes, no special SQLite build, and no
coordination protocol. It reads the WAL file directly from the filesystem.

\subsection{Beeper's Encryption Fork}

Beeper (acquired by Automattic, 2024) forked Litestream to add age
encryption for their Matrix bridge infrastructure. Their fork encrypts each
WAL segment and snapshot with an age recipient before uploading to S3. This
ensures that the S3 operator (or any party with bucket access) cannot read
the replicated data without the age identity (private key).

Beeper's fork demonstrated the viability of encrypted Litestream replication
in production at scale (millions of Matrix bridge databases). However, it
was never upstreamed and has diverged from mainline Litestream.

\subsection{age Encryption}

age~\cite{age} is a file encryption tool by Filippo Valsorda designed as a
modern replacement for PGP. It uses X25519 for key agreement, ChaCha20-Poly1305
for authenticated encryption, and HKDF for key derivation. age has three
properties that make it suitable for this use case:

\begin{enumerate}[nosep]
    \item \textbf{Small keys.} An age identity (private key) is a single
    62-character string. A recipient (public key) is a single 59-character
    string. Both fit in a Kubernetes Secret without base64 bloat.

    \item \textbf{Composable recipients.} A file can be encrypted to multiple
    recipients. This supports key rotation (encrypt to both old and new keys)
    and escrow (encrypt to both service key and recovery key).

    \item \textbf{No configuration.} age has no cipher suites, no key
    formats, no modes. There is one algorithm. There is one way.
\end{enumerate}

% ============================================================================
% 3. SPECIFICATION
% ============================================================================
